%This chapter was modified on 4/2/97.
%\setcounter{chapter}{1}
\chapter{Densitas Peluang Kontinu}\label{chp 2}

\section{Simulasi Peluang Kontinu}\label{sec 2.1}

Dalam bagian ini kita akan menunjukkan cara menggunakan simulasi komputer untuk
percobaan yang hasil-hasil mungkinnya membentuk suatu kontinuum\index{kontinuum}
penuh.

\subsection*{Peluang}

\begin{example}\label{exam 2.1.1}
Kita mulai dengan membuat sebuah pemutar\index{pemutar}, yang terdiri atas sebuah
lingkaran dengan \emx {keliling satu satuan} dan sebuah penunjuk seperti pada
Gambar~\ref{fig 2.05}. Kita pilih sebuah titik pada lingkaran itu dan memberinya
label 0, lalu memberi setiap titik lain pada lingkaran label berupa jarak,
misalnya $x$, dari 0 ke titik tersebut, yang diukur berlawanan arah jarum jam.
Percobaannya adalah memutar penunjuk dan mencatat label titik yang berada pada
ujung penunjuk. Kita gunakan peubah acak $X$ untuk menyatakan nilai hasil ini.
Ruang sampelnya jelas merupakan interval $[0, 1)$. Kita ingin membangun model
peluang yang membuat setiap hasil memiliki kemungkinan yang sama untuk terjadi.
\par
Jika kita melanjutkan seperti dalam Bab~\ref{chp 1} untuk percobaan dengan
sejumlah terhingga hasil yang mungkin, kita harus menetapkan peluang 0 bagi
setiap hasil. Jika tidak, jumlah peluang untuk seluruh hasil yang mungkin tidak
akan sama dengan 1. (Sebenarnya, menjumlahkan tak terhitung banyaknya bilangan
real merupakan perkara yang rumit; khususnya, agar penjumlahan semacam itu
bermakna, paling banyak hanya terhitung banyaknya suku yang boleh berbeda dari
0.) Namun, jika semua peluang yang ditetapkan adalah 0, jumlahnya menjadi 0,
bukan 1 sebagaimana seharusnya.
\putfig{2truein}{PSfig2-12arc}{Sebuah pemutar.}{fig 2.05}
\par
Pada bagian berikutnya, kita akan menunjukkan cara membangun model peluang dalam
situasi ini. Untuk sementara, kita akan menganggap bahwa model semacam itu dapat
dibangun. Kita juga akan menganggap bahwa dalam model tersebut, jika $E$ adalah
suatu busur lingkaran dan $E$ memiliki panjang $p$, maka model itu menetapkan peluang $p$
bagi $E$. Artinya, jika penunjuk diputar, peluang bahwa penunjuk akhirnya
mengarah ke suatu titik dalam $E$ sama dengan $p$, suatu hal yang tentu wajar
untuk diharapkan.
\par
Mensimulasikan percobaan ini dengan komputer merupakan hal yang mudah. Banyak
paket perangkat lunak komputer mempunyai fungsi yang menghasilkan bilangan real
acak dalam interval $[0, 1]$. Sebenarnya, nilai yang dihasilkan selalu berupa
bilangan rasional, dan nilai-nilai itu ditentukan oleh suatu algoritma, sehingga
barisan nilai tersebut tidak benar-benar acak. Meskipun demikian, barisan yang
dihasilkan algoritma semacam itu berperilaku sangat mirip dengan barisan acak
teoretis, sehingga dapat kita gunakan dalam simulasi percobaan. Sesekali kita
perlu merujuk fungsi semacam itu. Fungsi ini akan kita namai $rnd$\index{rnd}.
\end{example}

\subsection*{Prosedur Monte Carlo dan Luas}

Kadang-kadang kita perlu mengestimasi besaran yang nilai tepatnya sulit atau
mustahil dihitung secara eksak. Dalam beberapa kasus semacam itu, suatu prosedur
yang melibatkan keacakan, yang disebut \emx {prosedur Monte Carlo}, dapat
digunakan untuk memperoleh estimasi tersebut.

\begin{example}\label{exam 2.1.2}
Dalam contoh ini kita menunjukkan bagaimana simulasi dapat digunakan untuk
mengestimasi luas\index{luas, estimasi} bangun datar. Misalkan kita memprogram
komputer agar menghasilkan sepasang bilangan $(x,y)$, yang masing-masing dipilih
secara acak dan saling bebas dari interval $[0,1]$. Pasangan $(x,y)$ ini dapat
kita tafsirkan sebagai koordinat suatu titik yang dipilih \emx {secara acak}
dari persegi satuan. Kejadian adalah himpunan bagian dari persegi satuan.
Pengalaman kita pada Contoh~\ref{exam 2.1.1} menunjukkan bahwa titik tersebut
memiliki kemungkinan yang sama untuk jatuh pada himpunan-himpunan bagian yang
luasnya sama. Karena luas total persegi adalah 1, peluang bahwa titik jatuh pada
suatu himpunan bagian tertentu $E$ dari persegi satuan seharusnya sama dengan
luasnya. Dengan demikian, kita dapat mengestimasi luas setiap himpunan bagian
dari persegi satuan dengan mengestimasi peluang bahwa suatu titik yang dipilih
secara acak dari persegi itu jatuh dalam himpunan bagian tersebut.
\par
Kita dapat menggunakan metode ini untuk mengestimasi luas daerah $E$ di bawah
kurva $y = x^2$ dalam persegi satuan (lihat Gambar~\ref{fig 2.3}). Kita memilih
sejumlah besar titik $(x,y)$ secara acak dan mencatat proporsi titik yang jatuh
dalam daerah $E
= \{\,(x,y):y \leq x^2\,\}$.
\par
Program {\bf MonteCarlo}\index{MonteCarlo (program)} akan menjalankan percobaan
ini untuk kita. Menjalankan program ini untuk 10{,}000 percobaan menghasilkan
estimasi sebesar .325 (lihat Gambar~\ref{fig 2.4}).

\putfig{3truein}{PSfig2-3}{Luas di bawah $y = x^2.$}{fig 2.3}

Dari percobaan-percobaan ini kita akan mengestimasi luasnya sekitar 1/3. Tentu
saja, untuk daerah sederhana ini kita dapat mencari luas eksaknya dengan
kalkulus. Sesungguhnya,
$$
\mbox{Luas}\ E = \int_0^1 x^2\,dx = \frac13\ .
$$
Dalam Bab~\ref{chp 1} telah kita nyatakan bahwa ketika kita menyimulasikan
percobaan jenis ini sebanyak $n$ kali untuk mengestimasi suatu peluang, kita
dapat mengharapkan galat jawaban paling besar $1/\sqrt n$ dalam sekurang-
kurangnya 95 persen kasus. Untuk 10{,}000 percobaan kita dapat mengharapkan
ketelitian 0.01, dan simulasi kita memang mencapai ketelitian ini.

\putfig{3truein}{PSfig2-4}{Menghitung luas melalui simulasi.}{fig 2.4}  

Argumen yang sama berlaku untuk setiap daerah $E$ dalam persegi satuan. Sebagai
contoh, misalkan $E$ adalah lingkaran berpusat di $(1/2,1/2)$ dan berjari-jari
1/2. Peluang bahwa titik acak $(x,y)$ kita berada di dalam lingkaran sama dengan
luas lingkaran tersebut, yaitu
$$
P(E) = \pi{\Bigl(\frac{1}{2}\Bigr)}^2 = \frac{\pi}{4}\ .
$$
Jika kita tidak mengetahui nilai $\pi$, kita dapat mengestimasi
\index{$\pi$, estimasi|(} nilainya dengan melakukan percobaan ini berkali-kali!
\end{example}

Contoh di atas bukanlah satu-satunya cara mengestimasi nilai $\pi$ melalui
percobaan acak. Berikut ini cara lain yang ditemukan oleh Buffon.\footnote{G. L.
Buffon, in ``Essai d'Arithm\'etique Morale," \emx {Oeuvres Compl\`etes de
Buffon avec Supplements,} tome~iv, ed. Dum\'enil (Paris, 1836).}

\subsection*{Jarum Buffon}\index{BUFFON, G. L.}
\index{jarum Buffon|(}

\begin{example}\label{exam 2.1.3}
Misalkan kita mengambil sebuah meja kartu dan menggambar pada permukaan atasnya
sekumpulan garis sejajar yang terpisah sejauh satu satuan. Kemudian kita
menjatuhkan secara acak sebuah jarum biasa yang panjangnya satu satuan pada
permukaan itu dan mengamati apakah jarum tersebut melintasi salah satu garis.
Kita dapat menyatakan hasil-hasil yang mungkin dari percobaan ini dengan
koordinat sebagai berikut. Misalkan $d$ adalah jarak dari pusat jarum ke garis
terdekat. Selanjutnya, misalkan $L$ adalah garis yang ditentukan oleh jarum, dan
definisikan $\theta$ sebagai sudut lancip yang dibentuk garis $L$ dengan
kumpulan garis sejajar tersebut. (Pembaca tentu perlu berhati-hati terhadap
deskripsi ruang sampel ini. Kita sedang berusaha memberi koordinat kepada suatu
himpunan ruas garis. Untuk memahami mengapa pemilihan koordinat harus dilakukan
dengan cermat, lihat Contoh~\ref {exam 2.1.5}.) Dengan deskripsi ini, kita memiliki $0
\leq d \leq 1/2$, dan $0 \leq
\theta \leq \pi/2$. Selain itu, kita melihat bahwa jarum melintasi garis
terdekat jika dan hanya jika sisi miring segitiga (lihat Gambar~\ref
{fig 2.6}) lebih kecil daripada setengah panjang jarum, yaitu
$$
\frac d{\sin\theta} < \frac12\ .
$$

\putfig{4.5truein}{PSfig2-6}{Percobaan Buffon.}{fig 2.6}

Sekarang kita menganggap bahwa ketika jarum jatuh, pasangan $(\theta,d)$ dipilih
secara acak dari persegi panjang $0 \leq \theta \leq \pi/2$, $0 \leq d \leq
1/2$. Kita mengamati apakah jarum melintasi garis terdekat (yakni, apakah $d \leq
(1/2)\sin\theta$). Peluang kejadian $E$ ini adalah proporsi luas persegi panjang
yang berada di dalam $E$ (lihat Gambar~\ref{fig 2.7}). Luas persegi panjang
tersebut adalah $\pi/4$, sedangkan luas $E$ adalah
$$
\mbox{Luas} = \int_0^{\pi/2}\frac12\sin\theta\,d\theta = \frac 12\ .
$$
Karena itu, kita memperoleh
$$
P(E) = \frac{1/2}{\pi/4} = \frac2\pi\ .
$$

\putfig{3.5truein}{PSfig2-7}
{Himpunan $E$ dari pasangan $(\theta, d)$ dengan $d < {\frac{1}{2}} \sin \theta$.}{fig 2.7}

Program {\bf BuffonsNeedle}\index{BuffonsNeedle (program)} menyimulasikan
percobaan ini. Dalam Gambar~\ref{fig 2.8}, kita memperlihatkan posisi setiap
jarum ke-100 dalam satu kali pengoperasian program ketika 10{,}000 jarum
``dijatuhkan." Estimasi akhir kita untuk $\pi$ adalah 3.139. Walaupun hasil ini
berjarak kurang dari 0.003 terhadap nilai sejati $\pi$, kita tidak berhak
mengharapkan ketelitian sebesar itu. Alasannya adalah simulasi kita mengestimasi
$P(E)$. Meskipun kita dapat mengharapkan galat estimasi ini paling besar 0.001,
galat kecil dalam $P(E)$ menjadi lebih besar ketika kita menggunakannya untuk
menghitung $\pi = 2/P(E)$. Perlman\index{PERLMAN, M. D.} dan
Wichura\index{WICHURA, M. J.}, dalam artikel mereka
``Sharpening Buffon's Needle,"\footnote{M. D. Perlman and M. J. Wichura, ``Sharpening
Buffon's Needle," \emx {The American Statistician,} vol.~29, no.~4 (1975), pp.~157--163.}
menunjukkan bahwa sekitar 95 persen dari seluruh kesempatan, kita dapat
mengharapkan galat yang tidak melebihi $5/\sqrt n$. Di sini $n$ adalah jumlah
jarum yang dijatuhkan. Jadi, untuk 10{,}000 jarum kita seharusnya mengharapkan
galat tidak lebih dari 0.05, dan memang demikian dalam kasus ini. Kita melihat
bahwa diperlukan banyak percobaan untuk memperoleh estimasi $\pi$ yang cukup
baik.\index{$\pi$, estimasi|)}
\index{jarum Buffon|)}
\end{example}

\putfig{4truein}{PSfig2-8}{Simulasi percobaan jarum Buffon.}{fig 2.8}

Dalam setiap contoh kita sejauh ini, kejadian yang berukuran sama mempunyai
peluang yang sama. Berikut ini contoh ketika hal tersebut tidak berlaku. Nanti
kita akan melihat banyak contoh lain semacam ini.
\begin{example}\label{exam 2.1.4.5}
Misalkan kita memilih dua bilangan real acak dalam $[0,1]$ dan menjumlahkannya.
Misalkan $X$ adalah jumlah tersebut. Bagaimana distribusi $X$?
\par
Untuk membantu memahami jawaban atas pertanyaan ini, kita dapat menggunakan
program {\bf Areabargraph}\index{Areabargraph (program)}. Program ini
menghasilkan diagram batang dengan sifat bahwa pada setiap interval,
\emx {luas} batang, bukan tingginya, sama dengan proporsi hasil yang jatuh dalam
interval bersangkutan. Kita telah melakukan percobaan ini 1000 kali; datanya
ditampilkan dalam Gambar~\ref{fig 2.8.5}. Tampaknya fungsi yang didefinisikan
oleh
$$f(x) = \left \{ \begin{array}{ll}
                  x,   & \mbox{jika $0 \le x \le 1$,}  \\
                  2-x, & \mbox{jika $1 < x \le 2$} 
                  \end{array}
         \right.
$$
menyesuaikan diri dengan data dengan sangat baik. (Fungsi itu ditampilkan dalam
gambar.) Pada bagian berikutnya kita akan melihat bahwa fungsi inilah fungsi yang
``tepat." Maksudnya, jika $a$ dan $b$ adalah sembarang dua bilangan real antara
$0$ dan $2$, dengan $a \le b$, kita dapat menggunakan fungsi ini untuk
menghitung peluang bahwa $a \le X \le b$. Untuk memahami cara melakukan
perhitungan ini, kita kembali memperhatikan Gambar~\ref{fig 2.8.5}. Berdasarkan
cara batang-batang itu dibentuk, jumlah luas batang yang bersesuaian dengan
interval $[a, b]$ menghampiri peluang bahwa $a \le X \le b$. Namun, jumlah luas
batang-batang tersebut juga menghampiri integral
$$\int_a^b f(x)\,dx\ .$$ Hal ini menunjukkan bahwa untuk percobaan dengan suatu
kontinuum hasil yang mungkin, jika kita menemukan fungsi dengan sifat di atas,
kita akan dapat menggunakannya untuk menghitung peluang. Pada bagian berikutnya,
kita akan menunjukkan cara menentukan fungsi
\linebreak[4]$f(x)$.
\putfig{3.5truein}{PSfig2-8-5}{Jumlah dua bilangan acak.}{fig 2.8.5}
\end{example}

\begin{example}\label{exam 2.1.4.6}
Misalkan kita memilih 100 bilangan acak dalam $[0, 1]$, dan misalkan $X$
menyatakan jumlahnya. Bagaimana distribusi $X$? Kita telah melakukan percobaan
ini 10000 kali; hasilnya ditampilkan dalam Gambar~\ref{fig 2.8.6}. Dalam kasus
ini, fungsi apa yang menyesuaikan diri dengan batang-batang itu tidak begitu
jelas. Ternyata jenis fungsi yang dapat melakukannya disebut fungsi
\emx {densitas normal}\index{densitas normal}. Fungsi jenis ini kadang-kadang
disebut kurva ``berbentuk lonceng"\index{berbentuk lonceng}. Fungsi ini termasuk
di antara fungsi terpenting dalam bidang peluang dan akan didefinisikan secara
formal dalam Bagian~\ref{sec 5.2} dari Bab~\ref{chp 5}.
\putfig{3.5truein}{PSfig2-8-6}{Jumlah 100 bilangan acak.}{fig 2.8.6}
\end{example}
\par
Contoh terakhir kita menyelidiki pertanyaan mendasar tentang bagaimana peluang
ditetapkan.

\subsection*{Paradoks Bertrand}\index{paradoks Bertrand|(}

\begin{example}\label{exam 2.1.5} 
Tali busur suatu lingkaran adalah ruas garis yang kedua titik ujungnya terletak
pada lingkaran. Misalkan suatu tali busur dibuat \emx {secara acak}
\index{tali busur, acak} dalam lingkaran satuan. Berapakah peluang bahwa
panjangnya melebihi $\sqrt 3$?

Jawaban kita akan bergantung pada apa yang kita maksud dengan \emx {acak,} yang
pada gilirannya bergantung pada koordinat yang kita pilih. Ruang sampel
$\Omega$ adalah himpunan semua tali busur yang mungkin dalam lingkaran. Untuk
menentukan koordinat bagi tali-tali busur ini, mula-mula kita memperkenalkan
sistem koordinat persegi panjang dengan titik asal di pusat lingkaran (lihat
Gambar~\ref{fig 2.10}). Perhatikan bahwa suatu tali busur lingkaran tegak lurus
terhadap garis radial yang memuat titik tengah tali busur. Kita dapat
menyatakan setiap tali busur dengan memberikan:
 
\begin{enumerate}
\item koordinat persegi panjang $(x,y)$ dari titik tengah $M$, atau
\item koordinat polar $(r,\theta)$ dari titik tengah $M$, atau
\item koordinat polar $(1,\alpha)$ dan $(1,\beta)$ dari titik-titik ujung $A$
dan $B$.
 
\end{enumerate}
Dalam setiap kasus, kita akan menafsirkan \emx {secara acak} sebagai: memilih
koordinat-koordinat ini secara acak.

Kita dapat dengan mudah mengestimasi peluang ini melalui simulasi komputer.
Ketika memprogram simulasi ini, beberapa penyederhanaan akan memudahkan
pekerjaan; penyederhanaan tersebut kita uraikan satu per satu:

\putfig{2.5truein}{PSfig2-10}{Tali busur acak.}{fig 2.10}

\begin{enumerate}
\item Untuk menyimulasikan kasus ini, kita memilih secara acak nilai $x$ dan
$y$ dari $[-1,1]$. Kemudian kita memeriksa apakah $x^2 + y^2 \leq 1$. Jika
tidak, titik $M = (x,y)$ terletak di luar lingkaran sehingga tidak mungkin
menjadi titik tengah tali busur mana pun, dan kita mengabaikannya. Jika ya,
$M$ terletak di dalam lingkaran dan merupakan titik tengah suatu tali busur
tunggal, yang panjangnya $L$ diberikan oleh rumus:
$$
L = 2\sqrt{1 - (x^2 + y^2)}\ .
$$
\item Untuk menyimulasikan kasus ini, kita memanfaatkan fakta bahwa sebarang
rotasi lingkaran tidak mengubah panjang tali busur, sehingga sejak awal kita
dapat saja menganggap tali busurnya horizontal. Kemudian kita memilih $r$
secara acak dari $[-1,1]$ dan menghitung panjang tali busur yang dihasilkan,
dengan titik tengah $(r,\pi/2)$, menggunakan rumus:
$$
L = 2\sqrt{1 - r^2}\ .
$$
\item Untuk menyimulasikan kasus ini, kita menganggap salah satu titik ujung,
misalnya $B$, terletak di $(1, 0)$ (yakni, $\beta = 0$). Kemudian kita memilih
secara acak suatu nilai $\alpha$ dari $[0,2\pi]$ dan menghitung panjang tali
busur yang dihasilkan dengan menggunakan Hukum Kosinus melalui rumus:
$$
L = \sqrt{2 - 2\cos\alpha}\ .
$$
\end{enumerate}

Program {\bf BertrandsParadox}\index{BertrandsParadox (program)} menjalankan
simulasi ini. Menjalankan program tersebut menghasilkan hasil yang ditampilkan
dalam Gambar~\ref{fig 2.11}. Dalam lingkaran pertama pada gambar ini, telah
digambar sebuah lingkaran yang lebih kecil. Tali-tali busur yang memotong
lingkaran kecil ini mempunyai panjang sekurang-kurangnya $\sqrt 3$. Dalam
lingkaran kedua pada gambar, garis vertikal memotong semua tali busur yang
panjangnya sekurang-kurangnya $\sqrt 3$. Dalam lingkaran ketiga, garis vertikal
kembali memotong semua tali busur yang panjangnya sekurang-kurangnya $\sqrt 3$.

Dalam setiap kasus, kita menjalankan percobaan berkali-kali dan mencatat
proporsi panjang yang melebihi $\sqrt 3$. Kita telah mencetak hasil setiap
percobaan ke-100 sampai dengan 10{,}000 percobaan.

\putfig{4.5truein}{PSfig2-11}{Paradoks Bertrand.}{fig 2.11}

Menarik untuk mengamati bahwa ketiga proporsi ini \emx {tidak} sama; nilainya
bergantung pada koordinat yang kita pilih. Fenomena ini pertama kali diamati
oleh Bertrand dan kini dikenal sebagai \emx {paradoks Bertrand.}
\index{BERTRAND, J.}\footnote{J. Bertrand, \emx {Calcul des Probabilit\'es}
(Paris: Gauthier-Villars, 1889).} Sebenarnya ini sama sekali bukan paradoks;
fenomena ini hanya mencerminkan fakta bahwa pilihan koordinat yang berbeda akan
menghasilkan penetapan peluang yang berbeda. Penetapan mana yang ``benar"
bergantung pada penerapan atau penafsiran model yang dimaksud.
\par
Kita dapat membayangkan suatu percobaan nyata berupa melempar sedotan panjang
ke sebuah lingkaran yang digambar di atas meja kartu. Penetapan koordinat yang
``benar" seharusnya tidak bergantung pada letak lingkaran di atas meja kartu
atau letak meja kartu di dalam ruangan.
Jaynes\index{JAYNES, E. T.}\footnote{E. T. Jaynes, ``The Well-Posed Problem," in \emx
{Papers on Probability, Statistics and Statistical Physics,} R.~D.~Rosencrantz, ed.
(Dordrecht: D.~Reidel, 1983), pp.~133--148.} telah menunjukkan bahwa satu-satunya
penetapan yang memenuhi persyaratan ini adalah (2). Dalam pengertian ini,
penetapan (2) adalah penetapan yang alami, atau yang ``benar" (lihat
Latihan~\ref{exer
2.1.11}).                                             
\par
Dalam setiap kasus, kita dapat dengan mudah melihat peluang sejatinya jika
mengingat bahwa $\sqrt 3$ adalah panjang sisi segitiga sama sisi yang terinskripsi.
Karena itu, suatu tali busur mempunyai panjang $L > \sqrt 3$ jika titik
tengahnya berjarak $d < 1/2$ dari titik asal (lihat Gambar~\ref{fig 2.10}).
Perhitungan berikut menentukan peluang bahwa $L > \sqrt 3$ dalam masing-masing
dari ketiga kasus.
\begin{enumerate}
\item $L > \sqrt 3$ jika $(x,y)$ terletak di dalam lingkaran berjari-jari 1/2,
yang terjadi dengan peluang
$$
p = \frac{\pi(1/2)^2}{\pi(1)^2} = \frac14\ .
$$
\item $L > \sqrt 3$ jika $|r| < 1/2$, yang terjadi dengan peluang
$$
\frac{1/2 - (-1/2)}{1 - (-1)} = \frac12\ .
$$
\item $L > \sqrt 3$ jika $2\pi/3 < \alpha < 4\pi/3$, yang terjadi dengan peluang
$$
\frac{4\pi/3 - 2\pi/3}{2\pi - 0} = \frac13\ .
$$
\end{enumerate}
Kita melihat bahwa simulasi kita cukup sesuai dengan nilai-nilai teoretis ini.
\index{paradoks Bertrand|)}
\end{example}

\subsection*{Catatan Historis}

G.~L.~Buffon\index{BUFFON, G. L.|(} (1707--1788) adalah ilmuwan alam abad
kedelapan belas yang menerapkan peluang dalam sejumlah penyelidikannya. Karyanya termuat
dalam \emx {Histoire Naturelle} yang monumental, terdiri atas 44 jilid, beserta
suplemennya.\footnote{G.~L.~Buffon, \emx {Histoire Naturelle, Generali et
Particular avec le Descripti\'on du Cabinet du Roy,} 44~vols. (Paris:
L`Imprimerie Royale, 1749--1803).}
Sebagai contoh, ia menyajikan sejumlah tabel mortalitas dan menggunakannya untuk
menghitung harapan sisa usia bagi setiap kelompok umur. Dari tabelnya ia
mengamati bahwa harapan sisa usia seorang bayi berumur satu tahun adalah 33
tahun, sedangkan harapan sisa usia seorang pria berumur 21 tahun juga sekitar
33 tahun. Jadi, seorang ayah yang belum berumur 21 tahun dapat berharap hidup
lebih lama daripada anaknya yang berumur satu tahun. Namun, jika sang ayah
berumur 40 tahun, rasio peluang bahwa anaknya akan hidup lebih lama daripadanya
sudah 3 banding 2.\footnote{G.~L.~Buffon, ``Essai
d'Arithm\'etique Morale," p.~301.}
\par
Buffon ingin menunjukkan bahwa tidak semua perhitungan peluang hanya bertumpu
pada aljabar; sebagian bertumpu pada perhitungan geometris. Salah satunya adalah
``masalah jarum"\index{jarum Buffon|(} yang terkenal dan dibahas dalam bab
ini.\footnote{ibid., pp.~277--278.} Dalam rumusan aslinya, Buffon menggambarkan
suatu permainan ketika dua penjudi menjatuhkan sebongkah roti Prancis ke lantai
yang tersusun dari papan-papan lebar dan bertaruh apakah roti itu melintasi
celah lantai. Buffon bertanya: berapakah seharusnya panjang $L$ roti tersebut
dibandingkan dengan lebar $W$ papan lantai agar permainannya adil? Ia menemukan
jawaban yang benar ($L = (\pi/4)W$) dengan menggunakan metode yang pada dasarnya
sama dengan yang diuraikan dalam bab ini. Ia juga meninjau lantai berpola papan
catur, tetapi memberikan jawaban yang salah untuk kasus tersebut. Jawaban yang
benar kemudian diberikan oleh Laplace.\index{BUFFON, G. L.|)}\index{LAPLACE,
P. S.}
\par
Literatur memuat uraian sejumlah percobaan yang benar-benar dilakukan untuk
mengestimasi $\pi$ dengan metode menjatuhkan jarum ini.
N.~T.~Gridgeman\index{GRIDGEMAN, N. T.}\footnote{N.~T.~Gridgeman, ``Geometric Probability and
the Number $\pi$" \emx {Scripta Mathematika,} vol.~25, no.~3, (1960),
pp.~183--195.} membahas percobaan-percobaan yang ditampilkan dalam
Tabel~\ref{table 2.1}.
\begin{table}
\centering
$$ 
\begin{tabular}{llrll}
               &  Panjang                &  Jumlah          &  Jumlah         &  Estimasi  \\
Pelaku percobaan& jarum                  & lemparan\,\,\,\,\,\,&  perpotongan    &  $\pi$  \\ \hline
Wolf, 1850     &  \hspace{.073in}.8      & 5000\,\,\,\,\,\, &  2532                  & 3.1596      \\
Smith, 1855    &  \hspace{.073in}.6      & 3204\,\,\,\,\,\, &  1218.5                & 3.1553      \\
De Morgan, c.1860& 1.0                   & 600\,\,\,\,\,\,  &  \hspace{.075in}382.5  & 3.137       \\
Fox, 1864      &  \hspace{.078in}.75     & 1030\,\,\,\,\,\, &  \hspace{.075in}489    & 3.1595      \\
Lazzerini, 1901&  \hspace{.078in}.83     & 3408\,\,\,\,\,\, &  1808                  & 3.1415929   \\
Reina, 1925    &  \hspace{.083in}.5419   & 2520\,\,\,\,\,\, &  \hspace{.075in}869    & 3.1795  \\ \hline
\end{tabular}
$$
\caption{Percobaan jarum Buffon untuk mengestimasi $\pi$.}
\label{table 2.1}
\end{table}
(Bilangan setengah pada jumlah perpotongan berasal dari kompromi ketika tidak
dapat dipastikan apakah suatu perpotongan benar-benar terjadi.) Seperti kita,
ia mengamati bahwa 10{,}000 lemparan paling jauh hanya dapat menetapkan tempat
desimal pertama $\pi$ dengan tingkat keyakinan yang wajar. Gridgeman menunjukkan
bahwa walaupun tak satu pun percobaan menggunakan sampai 10{,}000 lemparan,
hasilnya ternyata sangat baik, bahkan dalam beberapa kasus terlalu baik.
Kenyataan bahwa jumlah lemparan tidak selalu berupa angka bulat yang rapi
menunjukkan bahwa para pelaku percobaan mungkin menggunakan penghentian
yang cerdik untuk memperoleh jawaban yang baik. Gridgeman berkomentar bahwa
estimasi Lazzerini ternyata sesuai dengan hampiran $\pi$ yang terkenal,
$355/113 = 3.1415929$, yang ditemukan oleh matematikawan Tiongkok abad kelima,
Tsu Ch'ungchih. Gridgeman mengatakan bahwa ia tidak memiliki laporan asli
Lazzerini. Sambil menunggunya (dan hanya mengetahui bahwa jarum memotong garis
1808 kali dalam 3408 lemparan), ia menyimpulkan bahwa panjang jarum pasti 5/6.
Ia menghitungnya dari rumus Buffon dengan menganggap $\pi = 355/113$:
$$
L = \frac{\pi P(E)}2 =
\frac12\left(\frac{355}{113}\right)\left(\frac{1808}{3408}\right) = \frac56 =
.8333\ .
$$
Bahkan dengan perencanaan yang cermat, seseorang harus sangat beruntung agar
dapat berhenti dengan begitu cerdik\index{jarum Buffon|)}.
\par
Penulis kedua gemar menelusuri awal minatnya pada teori peluang hingga Chicago
World's Fair\index{Chicago World's Fair} tahun 1933, tempat ia mengamati sebuah
alat mekanis yang menjatuhkan jarum dan menampilkan estimasi nilai $\pi$ yang
terus berubah. (Penulis pertama gemar menelusuri awal minatnya pada teori
peluang hingga penulis kedua.)

\exercises
\begin{LJSItem}

\istar\label{exer 2.1.1} Dalam masalah pemutar (lihat
Contoh~\ref{exam 2.1.1}), bagilah keliling satuan menjadi tiga busur dengan
panjang 1/2, 1/3, dan 1/6. Tulislah sebuah program untuk menyimulasikan
percobaan pemutar 1000 kali dan mencetak proporsi hasil yang jatuh dalam
masing-masing dari ketiga busur. Kemudian buatlah diagram batang yang
batang-batangnya memiliki lebar 1/2, 1/3, dan 1/6, serta luas yang sama dengan
proporsi terkait sebagaimana ditentukan oleh simulasi Anda. Tunjukkan bahwa
tinggi semua batang hampir sama.

\i\label{exer 2.1.2} Lakukan hal yang sama seperti dalam
Latihan~\ref{exer 2.1.1}, tetapi bagilah keliling satuan menjadi lima busur
dengan panjang 1/3, 1/4, 1/5, 1/6, dan 1/20.

\i\label{exer 2.1.3} Ubahlah program {\bf MonteCarlo} untuk mengestimasi luas
lingkaran berjari-jari 1/2 yang berpusat di $(1/2,1/2)$ di dalam persegi satuan
dengan memilih 1000 titik secara acak. Bandingkan hasil Anda dengan nilai sejati
$\pi/4$. Gunakan hasil tersebut untuk mengestimasi nilai $\pi$. Seberapa teliti
estimasi Anda?

\i\label{exer 2.1.4} Ubahlah program {\bf MonteCarlo} untuk mengestimasi luas
di bawah grafik $y = \sin\pi x$ di dalam persegi satuan dengan memilih
10{,}000 titik secara acak. Kemudian hitunglah nilai sejati luas ini dan
gunakan hasil Anda untuk mengestimasi nilai $\pi$. Seberapa teliti estimasi Anda?

\i\label{exer 2.1.5} Ubahlah program {\bf MonteCarlo} untuk mengestimasi luas
di bawah grafik $y = 1/(x + 1)$ dalam persegi satuan dengan cara yang sama
seperti dalam Latihan~\ref {exer 2.1.4}. Hitunglah nilai sejati luas ini dan
gunakan hasil simulasi Anda untuk mengestimasi nilai $\log 2$. Seberapa teliti
estimasi Anda?

\i\label{exer 2.1.6} Untuk menyimulasikan masalah jarum
Buffon\index{jarum Buffon}, kita memilih secara acak dan saling bebas jarak~$d$
serta sudut $\theta$, dengan $0 \leq d \leq 1/2$ dan $0 \leq \theta \leq
\pi/2$, lalu memeriksa apakah $d \leq (1/2)\sin\theta$. Dengan melakukannya
berkali-kali, kita mengestimasi $\pi$ sebagai $2/a$, dengan $a$ adalah proporsi
percobaan ketika $d \leq (1/2)\sin\theta$. Tulislah program untuk mengestimasi
$\pi$ dengan metode ini. Jalankan program Anda beberapa kali untuk masing-masing
100, 1000, dan 10{,}000 percobaan. Apakah ketelitian hampiran eksperimental
untuk $\pi$ meningkat seiring bertambahnya jumlah percobaan?

\i\label{exer 2.1.7} Untuk masalah jarum Buffon\index{jarum Buffon},
Laplace\index{LAPLACE,
P. S.}\footnote{P. S. Laplace, \emx {Th\'eorie Analytique des Probabilit\'es} (Paris: Courcier,
1812).} meninjau suatu kisi dengan garis \emx {horizontal} dan
\emx {vertikal} yang terpisah sejauh satu satuan. Ia menunjukkan bahwa peluang
sebuah jarum dengan panjang $L \leq 1$ memotong sekurang-kurangnya satu garis
adalah
$$
p = \frac{4L - L^2}\pi\ .
$$
Untuk menyimulasikan percobaan ini, kita memilih secara acak suatu sudut
$\theta$ antara 0~dan $\pi/2$, serta secara saling bebas dua bilangan $d_1$ dan
$d_2$ antara 0 dan $L/2$.
(Kedua bilangan itu menyatakan jarak dari pusat jarum ke garis horizontal dan
vertikal terdekat.) Jarum memotong suatu garis jika $d_1 \leq
(L/2)\sin\theta$ atau $d_2 \leq (L/2)\cos\theta$. Kita melakukannya berkali-kali
dan mengestimasi $\pi$ sebagai
$$
\bar \pi = \frac{4L - L^2}a\ ,
$$
dengan $a$ adalah proporsi percobaan ketika jarum memotong sekurang-kurangnya
satu garis. Tulislah program untuk mengestimasi $\pi$ dengan metode ini,
jalankan program Anda untuk 100, 1000, dan 10{,}000 percobaan, lalu bandingkan
hasil Anda dengan metode Buffon yang diuraikan dalam
Latihan~\ref{exer 2.1.6}. (Ambillah $L = 1$.)

\i\label{exer 2.1.8} Sebuah jarum panjang\index{jarum Buffon} dengan panjang
$L$ yang jauh lebih besar daripada 1 dijatuhkan pada kisi dengan garis
horizontal dan vertikal yang terpisah sejauh satu satuan. Kita akan melihat
(dalam Latihan~\ref{sec 6.3}.\ref{exer 6.3.29}) bahwa rata-rata jumlah garis
$a$ yang dipotong kira-kira
$$
a = \frac{4L}\pi\ .
$$
Untuk mengestimasi $\pi$ melalui simulasi, pilihlah secara acak suatu sudut
$\theta$ antara 0 dan $\pi/2$, lalu hitung $L\sin\theta + L\cos\theta$.
Besaran ini dapat digunakan sebagai jumlah garis yang dipotong. Ulangi proses
ini berkali-kali dan estimasikan $\pi$ dengan
$$
\bar \pi = \frac{4L}a\ ,
$$
dengan $a$ adalah rata-rata jumlah garis yang dipotong per percobaan. Tulislah
program untuk menyimulasikan percobaan ini dan jalankan program Anda untuk
jumlah percobaan sebesar 100, 1000, dan 10{,}000. Bandingkan hasil Anda dengan
metode Laplace atau Buffon untuk jumlah percobaan yang sama. (Gunakan $L = 100$.)

\medbreak\noindent
Latihan-latihan berikut melibatkan percobaan yang hasil-hasilnya tidak semuanya
memiliki peluang yang sama. Kita akan membahas percobaan semacam itu secara
terperinci pada bagian berikutnya, tetapi di sini Anda diajak menyelidiki
beberapa kasus sederhana.

\i\label{exer 2.1.9} Banyak masalah waktu tunggu mempunyai hasil dengan
\emx {distribusi eksponensial}\index{densitas eksponensial}
\index{fungsi densitas!eksponensial}. Kita akan melihat (dalam
Bagian~\ref{sec 5.2}) bahwa hasil semacam itu disimulasikan dengan menghitung
$(-1/\lambda)\log(\mbox{rnd})$, dengan $\lambda > 0$. Untuk waktu tunggu yang
dihasilkan dengan cara ini, rata-rata waktu tunggu adalah $1/\lambda$. Sebagai
contoh, waktu menunggu sebuah mobil melintas di jalan raya, atau waktu antara
emisi partikel dari sumber radioaktif, disimulasikan oleh suatu barisan bilangan
acak yang masing-masing dipilih dengan menghitung
$(-1/\lambda)\log(\mbox{rnd})$, dengan $1/\lambda$ sebagai rata-rata waktu
antara mobil atau emisi. Tulislah program untuk menyimulasikan waktu antarmobil
ketika rata-rata waktu antarmobil adalah 30 detik. Mintalah program Anda
menghitung diagram batang berdasarkan luas untuk waktu-waktu tersebut dengan membagi
interval waktu dari 0 sampai 120 menjadi 24 subinterval. Pada pasangan sumbu
yang sama, plotlah fungsi $f(x) = (1/30)e^{-(1/30)x}$. Apakah fungsi tersebut
sesuai dengan diagram batang itu?

\i\label{exer 2.1.10.5} Dalam Latihan~\ref{exer 2.1.9}, distribusinya
``muncul begitu saja." Dalam masalah ini, sekali lagi kita akan meninjau
percobaan yang hasil-hasilnya tidak memiliki peluang yang sama. Kita akan
menentukan suatu fungsi $f(x)$ yang dapat digunakan untuk menentukan peluang
kejadian-kejadian tertentu. Misalkan $T$ adalah segitiga siku-siku pada bidang
dengan titik-titik sudut $(0, 0),\ (1, 0),$ dan $(0,1)$. Percobaannya adalah
memilih suatu titik secara acak di bagian dalam $T$ dan hanya mencatat
koordinat-$x$ titik itu. Jadi, ruang sampelnya adalah himpunan $[0,1]$, tetapi
hasil-hasilnya tampaknya tidak memiliki peluang yang sama. Kita dapat
menyimulasikan percobaan ini dengan meminta komputer menghasilkan dua bilangan
real acak dalam $[0, 1]$, lalu mencatat bilangan pertama jika jumlah keduanya
kurang dari 1. Tulislah program ini dan jalankan untuk 10{,}000 percobaan.
Kemudian buatlah diagram batang hasilnya dengan membagi interval $[0, 1]$
menjadi 10 interval. Bandingkan diagram batang tersebut dengan fungsi
$f(x) = 2 - 2 x$. Selanjutnya, tunjukkan bahwa terdapat suatu konstanta $c$
sedemikian sehingga tinggi $T$ pada nilai koordinat-$x$ sebesar $x$ adalah
$c$ kali $f(x)$ untuk setiap $x$ dalam $[0, 1]$. Terakhir, tunjukkan bahwa
$$\int_0^1 f(x)\,dx = 1\ .$$
Bagaimana fungsi $f(x)$ dapat digunakan untuk menentukan peluang bahwa hasilnya
berada antara $.2$ dan $.5$?

\i\label{exer 2.1.11} Berikut ini cara lain memilih sebuah tali busur
\index{tali busur, acak} \emx {secara acak} pada lingkaran berjari-jari satu
satuan. Bayangkan kita mempunyai meja kartu yang panjang setiap sisinya 100.
Kita menempatkan sumbu-sumbu koordinat pada meja sedemikian sehingga setiap
sisi meja sejajar dengan salah satu sumbu dan pusat meja menjadi titik asal.
Kemudian kita menempatkan lingkaran berjari-jari satu satuan pada meja sehingga
pusat lingkaran menjadi titik asal. Sekarang pilihlah secara acak suatu titik
$(x_0,y_0)$ dalam persegi dan suatu sudut $\theta$ dalam interval
$(-\pi/2,\pi/2)$. Misalkan $m = \tan\theta$. Persamaan garis yang melalui
$(x_0,y_0)$ dengan kemiringan $m$ adalah
$$
y = y_0 + m(x - x_0)\ ,
$$
dan jarak garis ini dari pusat lingkaran (yakni, titik asal) adalah
$$
d = \left|\frac{y_0 - mx_0}{\sqrt{m^2 + 1}}\right|\ .
$$

Kita dapat menggunakan rumus jarak ini untuk memeriksa apakah garis memotong
lingkaran (yakni, apakah $d < 1$). Jika ya, tali busur yang dihasilkan kita
anggap sebagai tali busur {\em acak}. Ini menggambarkan percobaan menjatuhkan
sedotan panjang secara acak pada meja yang telah digambari lingkaran.

Tulislah program untuk menyimulasikan percobaan ini 10000 kali dan mengestimasi
peluang bahwa panjang tali busur lebih besar daripada $\sqrt3$. Bagaimana
perbandingan estimasi Anda dengan hasil Contoh~\ref{exam 2.1.5}?  
\end{LJSItem}


\section{Fungsi Densitas Kontinu}\label{sec 2.2}

Pada bagian sebelumnya kita telah melihat cara menyimulasikan percobaan dengan
suatu kontinuum penuh hasil yang mungkin dan telah memperoleh pengalaman dalam
memikirkan percobaan semacam itu. Sekarang kita beralih ke masalah umum
penetapan peluang bagi hasil dan kejadian dalam percobaan tersebut. Di sini
kita akan membatasi perhatian pada percobaan yang ruang sampelnya dapat
dipandang sebagai himpunan bagian yang dipilih secara sesuai dari garis,
bidang, atau ruang Euklides lainnya. Kita mulai dengan beberapa contoh
sederhana.

\subsection*{Pemutar}

\begin{example}\label{exam 2.2.1}
Percobaan pemutar\index{pemutar} yang diuraikan dalam
Contoh~\ref{exam 2.1.1} memiliki interval $[0, 1)$ sebagai himpunan hasil yang
mungkin. Kita ingin membangun model peluang yang membuat setiap hasil memiliki
peluang yang sama untuk terjadi. Kita telah melihat bahwa dalam model semacam
itu, setiap hasil harus diberi peluang 0. Ini sama sekali tidak berarti bahwa
peluang \emx {setiap} kejadian harus nol. Sebaliknya, jika peubah acak $X$
menyatakan hasilnya, maka peluang
$$P(\,0 \leq X \leq 1)$$
bahwa ujung penunjuk pemutar berhenti \emx {di suatu tempat} pada lingkaran
seharusnya sama dengan 1. Selain itu, peluang bahwa ujung penunjuk berhenti pada
setengah lingkaran bagian atas seharusnya sama dengan peluang untuk setengah
bagian bawah, sehingga
$$
P\biggl(0 \leq X < \frac12\biggr) = P\biggl(\frac12 \leq X <
1\biggr) = \frac12\ .
$$
Secara lebih umum, dalam model kita, kita menginginkan persamaan
$$
P(c \leq X < d) = d - c
$$
berlaku untuk setiap pilihan $c$ dan $d$.
\par
Jika kita menetapkan $E = [c, d]$, rumus di atas dapat ditulis dalam bentuk
$$P(E) = \int_E f(x)\,dx\ ,$$ dengan $f(x)$ adalah fungsi konstan bernilai 1.
Rumus ini seharusnya mengingatkan pembaca pada rumus yang bersesuaian dalam
kasus diskret untuk peluang suatu kejadian:
$$P(E) = \sum_{\omega \in E} m(\omega)\ .$$
Perbedaannya adalah bahwa dalam kasus kontinu, besaran yang diintegralkan,
$f(x)$, bukanlah peluang hasil $x$. (Namun, jika menggunakan infinitesimal,
kita dapat memandang $f(x)\,dx$ sebagai peluang hasil $x$.)
\par
Dalam kasus kontinu, kita akan menggunakan konvensi berikut. Jika himpunan hasil
merupakan himpunan bilangan real, setiap hasil akan dirujuk dengan huruf Romawi
kecil seperti $x$. Jika himpunan hasil merupakan himpunan bagian dari $R^2$,
setiap hasil akan dinyatakan dengan $(x, y)$. Dalam kedua kasus, kadang lebih
mudah merujuk suatu hasil dengan menggunakan $\omega$, seperti dalam
Bab \ref{chp 1}.
\putfig{3.5truein}{PSfig2-16-5}{Percobaan pemutar.}{fig 2.16.5}
\par
Gambar \ref{fig 2.16.5} memperlihatkan hasil 1000 putaran pemutar. Fungsi
$f(x)$ juga ditampilkan dalam gambar. Perhatikan bahwa luas di bawah $f(x)$ dan
di atas suatu interval tertentu kira-kira sama dengan proporsi hasil yang jatuh
dalam interval tersebut. Fungsi $f(x)$ disebut \emx {fungsi densitas}
\index{fungsi densitas} dari peubah acak $X$. Kenyataan bahwa luas di bawah
$f(x)$ dan di atas suatu interval bersesuaian dengan suatu peluang merupakan
sifat pendefinisi fungsi densitas. Definisi tepat fungsi densitas akan segera
diberikan.
\end{example}

\subsection*{Anak Panah}

\begin{example}\label{exam 2.2.2}
Permainan anak panah\index{anak panah} dilakukan dengan melempar anak panah ke
sebuah sasaran berbentuk lingkaran dengan \emx {jari-jari satu satuan.}
Misalkan kita melempar anak panah satu kali sehingga mengenai sasaran, lalu
mengamati tempat jatuhnya.

Untuk menyatakan hasil-hasil yang mungkin dari percobaan ini, wajar jika kita
mengambil himpunan $\Omega$ yang terdiri atas semua titik pada sasaran sebagai
ruang sampel. Titik-titik ini mudah dinyatakan dengan koordinat persegi panjang
terhadap sistem koordinat yang titik asalnya berada di pusat sasaran, sehingga
setiap pasangan koordinat $(x,y)$ dengan $x^2 + y^2 \leq 1$ menyatakan suatu
hasil percobaan yang mungkin. Dengan demikian, $\Omega = \{\,(x,y) : x^2 + y^2
\leq 1\,\}$ merupakan himpunan bagian bidang Euklides, dan kejadian $E =
\{\,(x,y) : y > 0\,\}$, misalnya, bersesuaian dengan pernyataan bahwa anak
panah jatuh pada setengah bagian atas sasaran, dan seterusnya. Kecuali terdapat
alasan untuk meyakini hal lain (dan untuk pemain ahli, alasan itu mungkin saja
ada!), wajar jika kita menganggap koordinat dipilih \emx {secara acak.} (Ketika
melakukannya dengan komputer, setiap koordinat dipilih secara seragam dari
interval $[-1, 1]$. Jika titik yang dihasilkan tidak terletak di dalam lingkaran
satuan, titik itu tidak dihitung.) Argumen yang digunakan dalam contoh
sebelumnya kemudian menunjukkan bahwa peluang setiap kejadian elementer, yang
hanya terdiri atas satu hasil, harus nol. Argumen itu juga menunjukkan bahwa
peluang kejadian berupa jatuhnya anak panah dalam sebarang himpunan bagian $E$
dari sasaran seharusnya ditentukan oleh proporsi luas sasaran yang berada dalam
$E$. Jadi,
$$
P(E) = \frac{\mbox{luas}\ E}{\mbox{luas\ sasaran}} = \frac{\mbox{luas}\ 
E}\pi\ .
$$
Ini dapat ditulis dalam bentuk
$$P(E) = \int_E f(x)\,dx\ ,$$
dengan $f(x)$ adalah fungsi konstan yang bernilai $1/\pi$. Secara khusus, jika
$E = \{\,(x,y) : x^2 + y^2 \leq a^2\,\}$ adalah kejadian bahwa anak panah
jatuh dalam jarak $a < 1$ dari pusat sasaran, maka
$$
P(E) = \frac{\pi a^2}\pi = a^2\ .
$$
Sebagai contoh, peluang bahwa anak panah berada dalam jarak 1/2 dari pusat
adalah 1/4.
\end{example}

\begin{example}\label{exam 2.2.3}
Dalam permainan anak panah\index{anak panah} yang ditinjau di atas, misalkan
kita tidak mengamati tempat anak panah jatuh, melainkan mengamati seberapa jauh
tempat jatuhnya dari pusat sasaran.
\par
Dalam kasus ini, kita mengambil himpunan $\Omega$ dari semua lingkaran yang
berpusat di pusat sasaran sebagai ruang sampel. Lingkaran-lingkaran ini mudah
dinyatakan dengan jari-jarinya, sehingga setiap lingkaran diidentifikasi oleh
jari-jarinya $r$, $0 \leq r \leq 1$. Dengan cara ini, kita dapat memandang
$\Omega$ sebagai himpunan bagian $[0,1]$ dari garis bilangan real.
\par
Peluang apa yang harus kita tetapkan bagi kejadian $E$ dari $\Omega$? Jika
$$E = \{\,r : 0 \leq r \leq a\,\}\ ,$$ 
maka $E$ terjadi jika anak panah jatuh dalam jarak $a$ dari pusat, yakni di
dalam lingkaran berjari-jari $a$. Dalam contoh sebelumnya kita telah melihat
bahwa berdasarkan asumsi kita, peluang kejadian ini diberikan oleh
$$
P([0,a]) = a^2\ .
$$
Secara lebih umum, jika
$$E = \{\,r : a \leq r \leq b\,\}\ ,$$ 
maka berdasarkan asumsi dasar kita,
\begin{eqnarray*}
P(E) = P([a,b]) & = & P([0,b]) - P([0,a]) \\
              & = & b^2 - a^2 \\
              & = & (b - a)(b + a) \\
              & = & 2(b - a)\frac{(b + a)}2\ .
\end{eqnarray*}

Jadi, $P(E) = $2(panjang $E$)(titik tengah $E$). Di sini kita melihat bahwa
peluang yang ditetapkan bagi interval $E$ tidak hanya bergantung pada
panjangnya, tetapi juga pada titik tengahnya (yakni, bukan hanya pada seberapa
panjang interval itu, melainkan juga pada letaknya). Secara kasar, dalam
percobaan ini, kejadian berbentuk $E = [a,b]$ lebih mungkin jika berada dekat
tepi sasaran dan kurang mungkin jika berada dekat pusat. (Pengalaman yang lazim
bagi pemula! Kesimpulannya mungkin berbeda jika pemula digantikan oleh seorang
ahli.)

Sekali lagi, kita dapat menyimulasikannya dengan komputer.
Kita membagi luas sasaran menjadi sepuluh daerah konsentris yang sama tebal.

\putfig{5truein}{PSfig2-15}{Distribusi jarak anak panah dalam 400 lemparan.}{fig 2.15}
\par
Program komputer {\bf Darts}\index{Darts (program)} melempar $n$ anak panah dan
mencatat proporsi keseluruhan yang jatuh dalam masing-masing daerah konsentris.
Program {\bf Areabargraph} kemudian membuat diagram batang dengan
\emx {luas} batang ke-$i$ sama dengan proporsi keseluruhan yang jatuh dalam
daerah ke-$i$. Menjalankan program untuk 1000 anak panah menghasilkan diagram
batang pada Gambar~\ref{fig 2.15}.
\par
Perhatikan bahwa di sini tinggi batang tidak semuanya sama, tetapi bertambah
kira-kira secara linear terhadap $r$. Bahkan, fungsi linear $y = 2r$ tampaknya
cukup sesuai dengan diagram batang kita. Ini menunjukkan bahwa peluang anak
panah jatuh dalam jarak $a$ dari pusat seharusnya diberikan oleh {\em luas} di
bawah grafik fungsi $y = 2r$ antara 0 dan $a$. Luas ini adalah $a^2$, yang
sesuai dengan peluang yang telah kita tetapkan di atas bagi kejadian tersebut.
\end{example}

\subsection*{Koordinat Ruang Sampel}

Contoh-contoh ini menunjukkan bahwa untuk percobaan kontinu semacam ini, kita
seharusnya menetapkan peluang bahwa hasil jatuh dalam suatu interval tertentu
melalui luas di bawah fungsi yang sesuai.
\par
Secara lebih umum, kita menganggap bahwa koordinat yang sesuai dapat
diperkenalkan pada ruang sampel $\Omega$, sehingga kita dapat memandang
$\Omega$ sebagai himpunan bagian dari ${\bf R}^n$. Ruang sampel semacam itu
kita sebut \emx {ruang sampel kontinu.}\index{ruang sampel!kontinu} Misalkan
$X$ adalah peubah acak yang menyatakan hasil percobaan. Peubah acak semacam itu
disebut \emx {peubah acak kontinu.}\index{peubah acak!kontinu} Selanjutnya kita
mendefinisikan fungsi densitas bagi $X$ sebagai berikut.

\subsection*{Fungsi Densitas Peubah Acak Kontinu}

\begin{definition}
Misalkan $X$ adalah peubah acak kontinu bernilai real. Suatu
\emx {fungsi densitas}\index{fungsi densitas} bagi $X$ adalah fungsi bernilai
real $f$ yang memenuhi
$$P(a \le X \le b) = \int_a^b f(x)\,dx$$
untuk semua $a,\ b \in {\bf R}$.
\end{definition}
Perhatikan bahwa \emx {tidak} semua peubah acak kontinu bernilai real memiliki
fungsi densitas. Namun, dalam buku ini kita hanya akan meninjau peubah acak
kontinu yang memiliki fungsi densitas.
\par
Dalam kaitannya dengan densitas $f(x)$, jika $E$ merupakan himpunan bagian dari
${\mat R}$, maka
$$
P(X \in E) = \int_Ef(x)\,dx\ .
$$
Notasi ini menganggap bahwa $E$ adalah himpunan bagian dari ${\mat R}$ yang
membuat $\int_E f(x)\,dx$ bermakna.

\begin{example}(Lanjutan Contoh~\ref{exam 2.2.1})\label{exam 2.2.5}
Dalam percobaan pemutar\index{pemutar}, kita memilih interval $0 \leq x < 1$
sebagai himpunan hasil dan memilih fungsi berikut sebagai fungsi densitas:
$$
f(x) =  \left \{ \begin{array}{ll}
                     1, & \mbox{jika $0 \leq x < 1$,} \\
                     0, & \mbox{selainnya.}
                      \end{array} 
             \right.
$$
Jika $E$ adalah kejadian bahwa ujung penunjuk pemutar jatuh pada setengah
lingkaran bagian atas, maka $E = \{\,x : 0 \leq x \leq 1/2\,\}$, sehingga
$$
P(E) = \int_0^{1/2} 1\,dx = \frac12\ .
$$
Secara lebih umum, jika $E$ adalah kejadian bahwa ujung penunjuk jatuh dalam
interval $[a,b]$, maka
$$
P(E) = \int_a^b 1\,dx = b - a\ .
$$
\end{example}

\begin{example}(Lanjutan Contoh~\ref{exam 2.2.2})
\label{exam 2.2.6}
Dalam percobaan permainan anak panah\index{anak panah} pertama, kita memilih
sebuah cakram berjari-jari satu satuan pada bidang sebagai ruang sampel dan
fungsi berikut sebagai fungsi densitas:
$$
f(x,y) = \left \{ \begin{array}{ll}
                1/\pi, & \mbox{jika $x^2 + y^2 \leq 1$,} \\
                0,     & \mbox{selainnya.}
                  \end{array}
         \right.
$$
Peluang bahwa anak panah jatuh di dalam himpunan bagian $E$ kemudian diberikan
oleh
\begin{eqnarray*}
P(E) & = & \int\,\int_E \frac1\pi \,dx\,dy\\
     & = & \frac1\pi \cdot (\mbox{luas}\,\,\,E)\ .
\end{eqnarray*}
\end{example}

Dalam kedua contoh ini, fungsi densitas bersifat konstan dan tidak bergantung
pada hasil tertentu. Percobaan yang koordinatnya dipilih \emx {secara acak}
sering kali dapat dinyatakan oleh fungsi densitas \emx {konstan}. Seperti dalam
Bagian~\ref{sec 1.2}, fungsi densitas semacam itu kita sebut \emx {seragam}
\index{fungsi densitas!seragam}\index{fungsi densitas seragam} atau
\emx {berpeluang sama.} Namun, tidak semua percobaan berjenis demikian.
  
\begin{example}(Lanjutan Contoh~\ref{exam 2.2.3})
\label{exam 2.2.7} 
Dalam percobaan permainan anak panah\index{anak panah} kedua, kita memilih
interval satuan pada garis bilangan real sebagai ruang sampel dan fungsi berikut
sebagai densitas:
$$f(r) = \left \{ \begin{array}{ll}
                 2r, & \mbox{jika  $0 < r < 1$,} \\
                 0,  & \mbox{selainnya.}
                  \end{array}
         \right.
$$
Peluang bahwa anak panah jatuh pada jarak $r$, $a \leq r \leq b$, dari pusat
sasaran kemudian diberikan oleh
\begin{eqnarray*}
P([a,b]) & = & \int_a^b 2r\,dr\\
       & = & b^2 - a^2\ .
\end{eqnarray*}
Sekali lagi, karena densitas kecil ketika $r$ dekat dengan 0 dan besar ketika
$r$ dekat dengan 1, kita melihat bahwa dalam percobaan ini anak panah lebih
mungkin jatuh dekat tepi sasaran daripada dekat pusat. Dalam kaitannya dengan
diagram batang Contoh~\ref{exam 2.2.3}, tinggi batang menghampiri fungsi
densitas, sedangkan luas batang menghampiri peluang subinterval (lihat
Gambar~\ref{fig 2.15}).
\end{example}
\par
Dalam contoh ini kita melihat bahwa, tidak seperti pada ruang sampel diskret,
nilai $f(x)$ dari fungsi densitas untuk hasil $x$ \emx {bukanlah} peluang
terjadinya $x$ (kita telah melihat bahwa peluang ini selalu 0), dan secara umum
$f(x)$ \emx {sama sekali bukan suatu peluang.} Dalam contoh ini, jika kita
mengambil $\lambda = 2$, maka $f(3/4) = 3/2$. Karena lebih besar daripada 1,
nilai ini tidak mungkin merupakan peluang.
\par
Meskipun demikian, fungsi densitas $f$ memang memuat seluruh informasi peluang
tentang percobaan karena peluang semua kejadian dapat diturunkan darinya.
Secara khusus, peluang bahwa hasil percobaan jatuh dalam interval $[a,b]$
diberikan oleh
$$
P([a,b]) = \int_a^b f(x)\,dx\ ,
$$
yakni oleh \emx {luas} di bawah grafik fungsi densitas pada interval $[a,b]$.
Jadi, di sini terdapat hubungan erat antara peluang dan luas. Hubungan erat ini
telah memandu kita dalam membuat diagram batang: setiap batang dipilih
sedemikian sehingga \emx {luasnya,} bukan tingginya, menyatakan frekuensi
relatif terjadinya hasil, dan dengan demikian mengestimasi peluang bahwa hasil
jatuh dalam interval terkait.
\par
Dalam bahasa kalkulus, kita dapat mengatakan bahwa peluang terjadinya kejadian
berbentuk $[x, x + dx]$, dengan $dx$ kecil, kira-kira diberikan oleh
$$
P([x, x+dx]) \approx f(x)dx\ ,
$$
yakni oleh luas persegi panjang di bawah grafik $f$. Perhatikan bahwa ketika
$dx \to 0$, peluang ini $\to 0$, sehingga peluang $P(\{x\})$ dari satu titik
kembali bernilai 0, seperti dalam Contoh~\ref{exam 2.2.1}.
\par
Sekilas pandang pada grafik fungsi densitas segera memberi tahu kita kejadian
mana dalam suatu percobaan yang lebih mungkin. Secara kasar, kita dapat
mengatakan bahwa di tempat densitasnya besar, kejadian lebih mungkin, sedangkan
di tempat densitasnya kecil, kejadian kurang mungkin. Dalam
Contoh~\ref{exam 2.1.4.5}, fungsi densitas mencapai nilai terbesar pada 1. Jadi,
untuk dua interval $[0, a]$ dan $[1, 1+a]$, dengan $a$ suatu bilangan real
positif kecil, kita melihat bahwa $X$ lebih mungkin mengambil nilai dalam
interval kedua daripada dalam interval pertama.

\subsection*{Fungsi Distribusi Kumulatif Peubah Acak\\ Kontinu}
Kita telah melihat bahwa fungsi densitas berguna ketika meninjau peubah acak
kontinu. Ada jenis fungsi lain yang berkaitan erat dengan fungsi densitas dan
juga sangat penting. Fungsi-fungsi ini disebut fungsi
\emx {distribusi kumulatif}\index{fungsi distribusi kumulatif}.
\begin{definition}
Misalkan $X$ adalah peubah acak kontinu bernilai real. Fungsi distribusi
kumulatif $X$ didefinisikan oleh persamaan
$$F_X(x) = P(X \le x)\ .$$
\end{definition}
Jika $X$ adalah peubah acak kontinu bernilai real yang memiliki fungsi
densitas, maka peubah tersebut juga memiliki fungsi distribusi kumulatif. Teorema berikut
menunjukkan bahwa kedua fungsi itu berkaitan dengan cara yang sangat baik.
\begin{theorem} Misalkan $X$ adalah peubah acak kontinu bernilai real dengan
fungsi densitas $f(x)$. Fungsi yang didefinisikan oleh
$$F(x) = \int_{-\infty}^x f(t)\,dt$$
merupakan fungsi distribusi kumulatif $X$. Selain itu, kita memiliki
$${{d\ }\over{dx}} F(x) = f(x)\ .$$
\proof Berdasarkan definisi,
$$F(x) = P(X \le x)\ .$$
Misalkan $E = (-\infty, x]$. Maka
$$P(X \le x) = P(X \in E)\ ,$$
yang sama dengan
$$\int_{-\infty}^x f(t)\,dt\ .$$
\par
Menerapkan Teorema Dasar Kalkulus pada persamaan pertama dalam pernyataan
teorema menghasilkan pernyataan kedua.
\end{theorem}
\par
Dalam banyak percobaan, fungsi densitas peubah acak yang bersangkutan mudah
dituliskan. Namun, sering kali fungsi distribusi kumulatif lebih mudah diperoleh
daripada fungsi densitas. (Tentu saja, setelah memperoleh fungsi distribusi
kumulatif, kita dapat dengan mudah memperoleh fungsi densitas melalui
diferensiasi, sebagaimana ditunjukkan teorema di atas.) Sekarang kita berikan
beberapa contoh yang memperlihatkan fenomena ini.
\begin{example}\label{exam 2.2.7.1}
Suatu bilangan real dipilih secara acak dari $[0, 1]$ dengan peluang seragam,
lalu bilangan itu dikuadratkan. Misalkan $X$ menyatakan hasilnya. Apa bentuk
fungsi distribusi kumulatif $X$? Apa bentuk densitas $X$?
\par
Kita mulai dengan menggunakan $U$ untuk menyatakan bilangan real yang dipilih.
Maka $X = U^2$. Jika $0 \le x \le 1$, kita memperoleh
\begin{eqnarray*}
F_X(x) & = & P(X \le x) \\
& = & P(U^2 \le x) \\
& = & P(U \le \sqrt x) \\
& = & \sqrt x\ .
\end{eqnarray*}
Jelas bahwa $X$ selalu mengambil nilai antara 0 dan 1, sehingga fungsi
distribusi kumulatif $X$ diberikan oleh
$$
F_X(x) = \left \{ \begin{array}{ll}
                                 0, & \mbox{jika $x \le 0$}, \\
                         {\sqrt x}, & \mbox{jika $0 \le x \le 1$}, \\
                                 1, & \mbox{jika $x \ge 1$}.
                    \end{array}
           \right.
$$
Dari sini kita dengan mudah menghitung bahwa fungsi densitas $X$ adalah
$$
f_X(x) = \left \{ \begin{array}{ll}
                                      0, & \mbox{jika $x \le 0$}, \\
                         1/(2{\sqrt x}), & \mbox{jika $0 \le x \le 1$}, \\
                                      0, & \mbox{jika  $x > 1$}.
                    \end{array}
           \right.
$$
Perhatikan bahwa $F_X(x)$ kontinu, tetapi $f_X(x)$ tidak. (Lihat
Gambar~\ref{fig 5.5}.)
\putfig{4.5truein}{PSfig5-5}{Distribusi dan densitas untuk $X = U^2$.}{fig 5.5} 
\end{example}

Ketika merujuk peubah acak kontinu $X$ (misalnya dengan fungsi densitas
seragam), lazim dikatakan bahwa ``$X$ \emx {berdistribusi} seragam pada
interval $[a, b]$." Fungsi distribusi kumulatif $X$ juga lazim disebut fungsi
distribusi $X$. Jadi, kata ``distribusi" digunakan dengan beberapa arti
berbeda dalam bidang peluang. (Ingat bahwa kata ini juga memiliki arti ketika
membahas peubah acak diskret.) Ketika merujuk fungsi distribusi kumulatif suatu
peubah acak kontinu $X$, kita akan selalu menggunakan kata ``kumulatif" sebagai
penjelas, kecuali penggunaan penjelas lain, seperti ``normal" atau
``eksponensial," sudah memperjelas maksudnya. Dalam bahasa Inggris, karena
frasa ``uniformly densitied on the interval $[a, b]$" tidak dapat diterima,
kita harus mengatakan ``uniformly distributed."


\begin{example}\label{exam 2.2.7.2}
Dalam Contoh~\ref{exam 2.1.4.5}, kita meninjau suatu peubah acak yang
didefinisikan sebagai jumlah\index{peubah acak seragam!jumlah dua peubah
kontinu} dari dua bilangan real acak yang dipilih secara seragam dari $[0, 1]$.
Misalkan peubah acak $X$ dan $Y$ menyatakan kedua bilangan real yang dipilih.
Definisikan $Z = X + Y$. Sekarang kita akan menurunkan bentuk fungsi distribusi
kumulatif dan fungsi densitas $Z$.

\putfig{3truein}{PSfig2-15-5}
{Perhitungan fungsi distribusi untuk Contoh \protect\ref{exam
2.2.7.2}\protect.}{fig 2.15.5} 


\putfig{4.5truein}{PSfig5-6}
{Fungsi distribusi dan densitas untuk Contoh \protect\ref{exam
2.2.7.2}\protect.}{fig 5.6} 

\par
Di sini, sebagai ruang sampel $\Omega$, kita mengambil persegi satuan dalam %${\rm{\bf R}}^2$
$\mat{R}^2$ dengan densitas seragam. Suatu titik
$\omega \in \Omega$ terdiri atas pasangan bilangan $(x, y)$ yang dipilih secara
acak. Maka $0 \leq Z\leq 2$. Misalkan $E_z$ menyatakan kejadian bahwa $Z \le
z$. Dalam Gambar~\ref{fig 2.15.5}, kita menampilkan himpunan $E_{.8}$.
Kejadian $E_z$, untuk setiap $z$ antara 0 dan 1, terlihat sangat mirip dengan himpunan
berarsir dalam gambar. Untuk $1 < z \le 2$, himpunan $E_z$ terlihat seperti
persegi satuan yang sudut kanan atasnya dipotong oleh sebuah segitiga. Sekarang
kita dapat menghitung distribusi peluang $F_Z$ dari $Z$; distribusi itu
diberikan oleh
\begin{eqnarray*}
F_Z(z) & = & P(Z \le z) \\
       & = & \mbox {Luas\ }E_z \\
       & = & \left \{\begin{array}{ll}
                               0, & \mbox{jika  $z < 0$}, \\
                        (1/2)z^2, & \mbox{jika  $0 \le z \le 1$}, \\
                1 - (1/2)(2-z)^2, & \mbox{jika  $1 \le z \le 2$}, \\
                               1, & \mbox{jika  $2 < z$}.
                     \end{array}
             \right.
\end{eqnarray*}
Fungsi densitas diperoleh dengan mendiferensiasikan fungsi ini:
$$
f_Z(z)  = \left \{\begin{array}{ll}
                        0, & \mbox{jika  $z < 0$}, \\
                        z, & \mbox{jika  $0 \le z \le 1$}, \\
                    2 - z, & \mbox{jika  $1 \le z \le 2$}, \\
                        0, & \mbox{jika  $2 < z$}.
                  \end{array}
          \right.
$$
Grafik fungsi-fungsi ini dapat dilihat pada Gambar~\ref{fig 5.6}.
\end{example}

\begin{example}\label{exam 2.2.7.3}
Dalam permainan anak panah\index{anak panah} yang diuraikan dalam
Contoh~\ref{exam 2.2.2}, %{exam 2.2.2}, 
apa bentuk distribusi jarak anak panah dari pusat sasaran? Apa bentuk densitasnya?

\putfig{2.5truein}{PSfig5-8}
{Perhitungan $F_{z}$ untuk Contoh~\protect\ref{exam 2.2.7.3}\protect.}{fig 5.8} 

Di sini, seperti sebelumnya, ruang sampel $\Omega$ adalah cakram satuan dalam %${\rm{\bf R}}^2$,
$\mat{R}^2$, dengan koordinat $(X, Y)$. Misalkan $Z = \sqrt{X^2 + Y^2}$
menyatakan jarak dari pusat sasaran. Misalkan $E$ adalah kejadian $\{Z \le
z\}$. Fungsi distribusi $F_Z$ dari $Z$ (lihat Gambar~\ref{fig 5.8}) kemudian
diberikan oleh
\begin{eqnarray*}
F_Z(z) & = & P(Z \le z) \\
       & = & {{\mbox {Luas\ }E}\over \mbox {Luas\ sasaran}}\ .
\end{eqnarray*}
Jadi, kita dengan mudah menghitung bahwa
$$
F_Z(z) = \left \{ \begin{array}{ll}
                           0, & \mbox{jika  $z \le 0$}, \\
                         z^2, & \mbox{jika  $0 \le z \le 1$}, \\
                           1, & \mbox{jika  $z > 1$}.
                    \end{array}
           \right.
$$
Densitas $f_Z(z)$ kembali diberikan oleh turunan $F_Z(z)$:
$$
f_Z(z) = \left \{ \begin{array}{ll}
                          0, & \mbox{jika  $z \le 0$}, \\
                         2z, & \mbox{jika  $0 \le z \le 1$}, \\
                          0, & \mbox{jika  $z > 1$}.
                    \end{array}
           \right.
$$
Grafik fungsi-fungsi ini dapat dilihat pada Gambar~\ref{fig 5.9}.
\par
Kita dapat memverifikasi hasil ini melalui simulasi sebagai berikut. Kita
memilih secara acak nilai $X$ dan $Y$ dari $[0,1]$ dengan distribusi seragam,
menghitung $Z = \sqrt{X^2 + Y^2}$, memeriksa apakah $0 \leq Z \leq 1$, lalu
menyajikan hasilnya dalam diagram batang (lihat Gambar~\ref{fig 5.10}).
\putfig{5truein}
{PSfig5-9}
{Distribusi dan densitas untuk $Z =\protect\sqrt{X^2 + Y^2}\protect$.}{fig 5.9}
\putfig{3.5truein}{PSfig5-10}
{Hasil simulasi untuk Contoh~\protect\ref{exam 2.2.7.3}\protect.}{fig 5.10} 
\end{example}

\begin{example}\label{exam 2.2.7.4} %{exam 5.2.6}
Misalkan Tn.\ dan Ny.\ Lockhorn\index{Lockhorn, Tn.\ dan Ny.} sepakat bertemu
di Hanover Inn\index{Hanover Inn} antara pukul 5:00 dan
6:00~{\footnotesize P.M.} pada hari Selasa. Misalkan masing-masing tiba pada
waktu antara 5:00 dan 6:00 yang dipilih secara acak dengan peluang seragam.
Apa bentuk fungsi distribusi untuk lamanya orang yang tiba lebih dahulu harus
menunggu yang lain? Apa bentuk fungsi densitasnya?
\par
Di sini kita kembali dapat menggunakan persegi satuan untuk menyatakan ruang
sampel, dan $(X, Y)$ sebagai waktu kedatangan keluarga Lockhorn (setelah
5:00~{\footnotesize P.M.}). Misalkan $Z = |X - Y|$. Maka kita memiliki
$F_X(x) = x$ dan $F_Y(y) = y$. Selain itu (lihat Gambar~\ref{fig 5.11}),
\begin{eqnarray*}
F_Z(z) & = & P(Z \leq z) \\
       & = & P(|X - Y| \leq z) \\
       & = & \mbox {Luas\ }E\ .
\end{eqnarray*}
Dengan demikian, kita memiliki
$$
F_Z(z) = \left \{ \begin{array}{ll}
                         0, & \mbox{jika  $z \le 0$}, \\
             1 - (1 - z)^2, & \mbox{jika  $0 \le z \le 1$}, \\
                         1, & \mbox{jika  $z > 1$}.
                    \end{array}
           \right.
$$
\putfig{3.5truein}{PSfig5-11}{Perhitungan $F_{Z}$.}{fig 5.11} %4truein 
Densitas $f_Z(z)$ kembali diperoleh melalui diferensiasi:
$$
f_Z(z) = \left \{ \begin{array}{ll}
                         0, & \mbox{jika  $z \le 0$}, \\
                    2(1-z), & \mbox{jika  $0 \le z \le 1$}, \\
                         0, & \mbox{jika  $z > 1$}.
                    \end{array}
           \right.
$$

\end{example}

\begin{example}\label{exam 2.2.7.5}
Dalam banyak keadaan kita mengamati suatu barisan peristiwa yang terjadi pada
waktu ``acak." Sebagai contoh, kita mungkin mengamati emisi isotop
radioaktif\index{isotop radioaktif}, mobil yang melewati patok jarak di jalan
raya\index{mobil di jalan raya}, atau bola lampu\index{bola lampu} yang putus.
Dalam kasus semacam itu, kita dapat mendefinisikan peubah acak $X$ untuk
menyatakan waktu antara dua peristiwa berturut-turut. Jelas bahwa $X$ adalah
peubah acak kontinu yang daerah nilainya terdiri atas bilangan real
non-negatif. Sering kali kita dapat memodelkan $X$ menggunakan
\emx {densitas eksponensial}\index{densitas eksponensial}
\index{fungsi densitas!eksponensial}. Densitas ini diberikan oleh rumus
$$f(t) = \left \{ \begin{array}{ll}
                         \lambda e^{-\lambda t}, & \mbox{jika  $t \ge 0$}, \\
                                              0, & \mbox{jika  $t < 0$}.
                    \end{array}
         \right. 
$$
Bilangan $\lambda$ adalah bilangan real non-negatif dan menyatakan kebalikan
dari nilai rata-rata $X$. (Hal ini akan ditunjukkan dalam
Bab~\ref{chp 6}.) Jadi, jika rata-rata waktu antara dua peristiwa adalah 30
menit, maka $\lambda = 1/30$. Grafik fungsi densitas ini dengan $\lambda =
1/30$ ditampilkan dalam Gambar~\ref{fig 2.16}.
\putfig{3.5truein}{PSfig2-16}{Densitas eksponensial dengan $\lambda = 1/30$.}{fig 2.16} 
Dari gambar terlihat bahwa walaupun nilai rata-ratanya 30, sesekali $X$
mengambil nilai yang jauh lebih besar.
\par
Misalkan kita telah membeli komputer yang berisi cakram keras Warp
9\index{cakram keras, Warp 9}. Penjualnya mengatakan bahwa rata-rata waktu
antara dua kerusakan cakram keras jenis ini adalah 30 bulan. Lamanya waktu
antara dua kerusakan sering dianggap berdistribusi menurut densitas
eksponensial. Kita akan menganggap model ini berlaku di sini, dengan
$\lambda = 1/30$.
\par
Sekarang misalkan kita telah mengoperasikan komputer selama 15 bulan. Kita
menganggap cakram keras aslinya masih berfungsi. Kita bertanya berapa lama lagi
cakram keras itu dapat diharapkan terus berfungsi. Wajar jika seseorang
mengharapkan cakram keras akan berfungsi rata-rata selama 15 bulan lagi.
(Seseorang juga mungkin menduga bahwa cakram itu akan berfungsi lebih dari 15
bulan karena fakta bahwa cakram tersebut telah berfungsi selama 15 bulan
menunjukkan bahwa produk kita bukan barang cacat.) Waktu yang masih harus kita
tunggu merupakan peubah acak baru yang akan kita sebut $Y$. Jelas, $Y = X -
15$. Kita dapat menulis program komputer untuk menghasilkan barisan nilai-$Y$
simulasi. Untuk melakukannya, mula-mula kita menghasilkan barisan nilai $X$ dan
membuang nilai yang kurang dari atau sama dengan 15 (nilai-nilai ini
bersesuaian dengan kasus ketika cakram keras berhenti berfungsi sebelum 15
bulan). Untuk menyimulasikan suatu nilai $X$, kita menghitung nilai ekspresi
$$\Bigl(-{1\over{\lambda}}\Bigr)\log(rnd)\ ,$$
dengan $rnd$ menyatakan bilangan real acak antara 0 dan 1. (Bahwa ekspresi ini
memiliki densitas eksponensial akan ditunjukkan dalam Bab \ref{chp 5}.)
Gambar \ref{fig 2.17} memperlihatkan diagram batang berdasarkan luas untuk
10{,}000 nilai-$Y$ simulasi.
\par
Nilai rata-rata $Y$ dalam simulasi ini adalah 29.74, yang lebih dekat dengan
rata-rata masa pakai semula sebesar 30 bulan daripada nilai 15 bulan yang
diduga di atas. Selain itu, distribusi $Y$ tampak dekat dengan distribusi $X$.
Sesungguhnya, $X$ dan $Y$ memiliki distribusi yang sama. Sifat ini disebut
\emx {sifat tanpa ingatan}\index{sifat tanpa ingatan}, karena lamanya waktu yang
harus kita tunggu hingga suatu peristiwa terjadi tidak bergantung pada berapa
lama kita telah menunggu. Satu-satunya fungsi densitas kontinu yang mempunyai
sifat ini adalah densitas eksponensial.
\putfig{3.5truein}{PSfig2-17}{Sisa masa pakai cakram keras.}{fig 2.17} 
\end{example}

\subsection*{Penetapan Peluang}

Pertanyaan mendasar dalam praktik adalah: bagaimana kita memilih fungsi
densitas peluang untuk menyatakan suatu percobaan tertentu? Jawabannya sangat
bergantung pada jumlah dan jenis informasi yang tersedia bagi kita tentang
percobaan tersebut. Dalam beberapa kasus, kita dapat melihat bahwa hasil-
hasilnya memiliki peluang yang sama. Dalam kasus lain, kita dapat melihat bahwa
percobaannya menyerupai percobaan lain yang telah dinyatakan oleh densitas yang
diketahui. Dalam kasus lainnya lagi, kita dapat menjalankan percobaan
berkali-kali dan membuat dugaan yang wajar tentang densitas berdasarkan
distribusi hasil yang diamati, seperti yang kita lakukan dalam
Bab~\ref{chp 1}. Secara umum, masalah memilih fungsi densitas yang tepat untuk
suatu percobaan merupakan masalah utama bagi pelaku percobaan dan tidak selalu
mudah diselesaikan (lihat Contoh~\ref{exam 2.1.5}). Di sini kita tidak akan
membahas pertanyaan ini secara terperinci, melainkan akan menganggap bahwa
densitas yang tepat sudah diketahui untuk setiap percobaan yang sedang ditinjau.

Pengenalan koordinat yang sesuai untuk menyatakan ruang sampel kontinu, serta
densitas yang sesuai untuk menyatakan peluangnya, tidak selalu begitu jelas,
sebagaimana ditunjukkan contoh terakhir kita.

\subsection*{Pohon Tak Hingga}

\begin{example}\label{exam 2.2.12}
Tinjaulah percobaan berupa pelemparan koin adil secara berulang tanpa berhenti.
Dalam Contoh~\ref{exam 1.5} telah kita lihat bahwa untuk koin yang dilempar
$n$ kali, ruang sampel alaminya adalah pohon biner dengan $n$ tahap. Berdasarkan
hal ini, kita mengharapkan bahwa untuk koin yang dilempar berulang kali, ruang
sampel alaminya adalah pohon biner\index{diagram pohon!biner tak hingga} dengan
tak hingga banyaknya tahap, seperti ditunjukkan dalam Gambar~\ref{fig 2.23}.

Mengejutkan bahwa walaupun pohon $n$ tahap jelas merupakan ruang sampel
terhingga, pohon tanpa batas dapat dinyatakan sebagai ruang sampel kontinu.
Untuk melihat bagaimana hal ini terjadi, kita sepakati bahwa suatu hasil khas
dari percobaan pelemparan koin tanpa batas dapat dinyatakan oleh barisan
berbentuk $\omega = \{\mbox{H H T H T T H}\dots\}$. Jika kita menulis 1 untuk H
dan 0 untuk T, maka $\omega = \{1\ 1\ 0\ 1\ 0\ 0\ 1\dots\}$. Dengan cara ini,
setiap hasil dinyatakan oleh barisan angka 0 dan 1.

\putfig{4truein}{PSfig2-23}{Pohon untuk tak hingga banyaknya pelemparan koin.}{fig 2.23}
\par
Sekarang misalkan kita memandang barisan angka 0 dan 1 ini sebagai ekspansi
biner suatu bilangan real $x = .1101001\cdots$ yang terletak antara 0 dan 1.
(Suatu \emx {ekspansi biner}\index{ekspansi biner} serupa dengan ekspansi
desimal, tetapi berlandaskan 2, bukan 10.) Setiap hasil kemudian dinyatakan oleh
suatu nilai $x$, dan dengan cara ini $x$ menjadi koordinat bagi ruang sampel,
yang mengambil semua nilai real antara 0 dan 1. (Perhatikan bahwa dua barisan
berbeda mungkin bersesuaian dengan bilangan real yang sama. Sebagai contoh,
barisan $\{\mbox{T H H H H H}\ldots\}$ dan
$\{\mbox{H T T T T T}\ldots\}$ sama-sama bersesuaian dengan bilangan real
$1/2$. Di sini kita tidak akan mempersoalkan masalah yang tampak ini.)
\par
Peluang apa yang harus ditetapkan bagi kejadian dalam ruang sampel ini?
Sebagai contoh, tinjaulah kejadian $E$ yang terdiri atas semua hasil ketika
lemparan pertama menghasilkan kepala dan lemparan kedua menghasilkan ekor.
Setiap hasil semacam itu berbentuk $.10****\cdots$, dengan $*$ dapat berupa 0
atau 1. Jika $x$ adalah koordinat bernilai real kita, nilai $x$ untuk setiap
hasil semacam itu harus terletak antara $1/2 = .10000\cdots$ dan $3/4 =
.11000\cdots$. Selain itu, setiap nilai $x$ antara 1/2 dan 3/4 memiliki ekspansi
biner berbentuk $.10****\cdots$. Artinya, $\omega\in E$ jika dan hanya jika
$1/2 \leq x < 3/4$, dan dengan demikian kita melihat bahwa $E$ dapat dinyatakan
oleh interval $[1/2,3/4)$. Secara lebih umum, setiap kejadian yang terdiri atas
hasil dengan hasil-hasil dari $n$ lemparan pertama yang telah ditentukan
dinyatakan oleh interval biner berbentuk $[k/2^n,(k+1)/2^n)$.
\par
Dalam Bagian~\ref{sec 1.2} telah kita lihat bahwa pada percobaan yang melibatkan
$n$ lemparan, peluang setiap satu hasil harus tepat $1/2^n$. Oleh karena itu,
dalam percobaan pelemparan tanpa batas, peluang setiap kejadian yang terdiri
atas hasil dengan hasil-hasil dari $n$ lemparan pertama yang telah ditentukan
juga harus $1/2^n$. Namun, $1/2^n$ tepat sama dengan panjang interval nilai
$x$ yang menyatakan $E$! Jadi, kita melihat bahwa seperti pada percobaan
pemutar, peluang suatu kejadian $E$ ditentukan oleh proporsi interval satuan
yang terletak dalam $E$.

Tinjau kembali pernyataan: peluang munculnya kepala ketika koin adil dilempar
adalah 1/2. Kita telah menunjukkan bahwa salah satu penafsiran pernyataan ini
adalah bahwa jika koin dilempar tanpa henti, proporsi kepala akan mendekati 1/2.
Artinya, dalam korespondensi kita dengan barisan biner, kita mengharapkan
memperoleh barisan biner dengan proporsi angka 1 yang menuju 1/2. Kejadian $E$
yang terdiri atas barisan biner dengan sifat ini merupakan himpunan bagian
sejati dari himpunan semua barisan biner yang mungkin. Sebagai contoh, kejadian
ini tidak memuat barisan $011011011\ldots$ (yakni, (011) yang diulang
terus-menerus). Kejadian $E$ sesungguhnya merupakan himpunan bagian barisan
biner yang sangat rumit, tetapi peluangnya dapat ditentukan sebagai limit
peluang kejadian dengan sejumlah terhingga hasil, yang peluangnya diberikan
oleh ukuran pohon terhingga. Ketika peluang $E$ dihitung dengan cara ini,
nilainya ternyata 1. Hasil yang luar biasa ini dikenal sebagai
\emx {Hukum Kuat Bilangan Besar}\index{Hukum Bilangan Besar!Kuat}
\index{Hukum Kuat Bilangan\\ Besar} (atau
\emx {Hukum Rata-Rata})\index{Hukum Rata-Rata} dan merupakan salah satu
pembenaran bagi konsep frekuensi peluang\index{konsep frekuensi peluang} kita.
Kita akan membuktikan bentuk lemah teorema ini dalam Bab~\ref{chp 8}.
\end{example}

\exercises
\begin{LJSItem}
 
\i\label{exer 2.2.1} Misalkan Anda memilih \emx {secara acak} suatu bilangan
real $X$ dari interval $[2,10]$.

\begin{enumerate}
\item Carilah fungsi densitas $f(x)$ dan peluang suatu kejadian $E$ untuk
percobaan ini, dengan $E$ merupakan subinterval $[a,b]$ dari $[2,10]$.

\item Dari (a), carilah peluang bahwa $X > 5$, bahwa $5 < X < 7$, dan bahwa
$X^2 - 12X + 35 > 0$.
\end{enumerate}

\i\label{exer 2.2.2} Misalkan Anda memilih suatu bilangan real $X$ dari
interval $[2,10]$ dengan fungsi densitas berbentuk
$$
f(x) = Cx\ ,
$$
dengan $C$ suatu konstanta.
\begin{enumerate}

\item Carilah $C$.

\item Carilah $P(E)$, dengan $E = [a,b]$ merupakan subinterval dari $[2,10]$.

\item Carilah $P(X > 5)$, $P(X < 7)$, dan $P(X^2 - 12X + 35 > 0)$.
\end{enumerate}

\i\label{exer 2.2.100} Sama seperti Latihan \ref{exer 2.2.2}, tetapi misalkan
$$
f(x) = \frac Cx\ .
$$

\i\label{exer 2.2.4} Misalkan Anda melempar anak panah\index{anak panah} ke
sasaran berbentuk lingkaran berjari-jari 10 inci. Dengan menganggap bahwa Anda
mengenai sasaran dan koordinat hasilnya dipilih secara acak, carilah peluang
bahwa anak panah jatuh

\begin{enumerate}
\item dalam jarak 2 inci dari pusat.

\item dalam jarak 2 inci dari tepi.

\item dalam kuadran pertama sasaran.

\item dalam kuadran pertama dan dalam jarak 2 inci dari tepi.
\end{enumerate}

\i\label{exer 2.2.5} Misalkan Anda sedang mengamati sumber
radioaktif\index{isotop radioaktif} yang memancarkan partikel dengan laju yang
dinyatakan oleh densitas eksponensial
$$
f(t) = \lambda e^{-\lambda t}\ ,
$$
dengan $\lambda = 1$, sehingga peluang $P(0,T)$ bahwa suatu partikel akan
muncul dalam $T$ detik berikutnya adalah
$P([0,T]) = \int_0^T\lambda e^{-\lambda t}\,dt$. Carilah peluang bahwa suatu partikel
(tidak harus yang pertama) akan muncul

\begin{enumerate}
\item dalam satu detik berikutnya.

\item dalam 3 detik berikutnya.

\item antara 3 dan 4 detik dari sekarang.

\item setelah 4 detik dari sekarang.
\end{enumerate}

\i\label{exer 2.2.6} Anggaplah sebuah bola lampu baru\index{bola lampu} akan
putus setelah $t$ jam, dengan $t$ dipilih dari $[0,\infty)$ menurut densitas
eksponensial
$$
f(t) = \lambda e^{-\lambda t}\ .
$$
Dalam konteks ini, $\lambda$ sering disebut \emx {laju kegagalan} bola lampu.

\begin{enumerate}
\item Anggaplah $\lambda = 0.01$, lalu carilah peluang bahwa bola lampu
\emx {tidak} akan putus sebelum $T$ jam. Peluang ini sering disebut
\emx {keandalan} bola lampu.

\item Untuk nilai $T$ berapakah keandalan bola lampu $ = 1/2$?
\end{enumerate}

\i\label{exer 2.2.7} Pilihlah suatu bilangan $B$ \emx {secara acak} dari
interval $[0,1]$ dengan densitas seragam. Carilah peluang bahwa
\begin{enumerate}
\item $1/3 < B < 2/3$.

\item $|B - 1/2| \leq 1/4$.

\item $B < 1/4$ atau $1 - B < 1/4$.

\item $3B^2 < B$.
\end{enumerate}

\i\label{exer 2.2.8} Pilihlah secara saling bebas dua bilangan $B$ dan $C$
\emx {secara acak} dari interval $[0,1]$ dengan densitas seragam. Perhatikan
bahwa titik $(B,C)$ kemudian dipilih \emx {secara acak} dalam persegi satuan.
Carilah peluang bahwa

\begin{enumerate}
\item $B + C < 1/2$.

\item $BC < 1/2$.

\item $|B - C| < 1/2$.

\item $\max\{B,C\} < 1/2$.

\item $\min\{B,C\} < 1/2$.

\item $B < 1/2$ dan $1 - C < 1/2$.

\item syarat (c) dan (f) keduanya berlaku.

\item $B^2 + C^2 \leq 1/2$.

\item $(B - 1/2)^2 + (C - 1/2)^2 < 1/4$.
\end{enumerate}

\i\label{exer 2.2.8.5} Misalkan kita mempunyai suatu barisan peristiwa. Kita
menganggap bahwa waktu $X$ antara dua peristiwa berdistribusi eksponensial
dengan $\lambda = 1/10$, sehingga rata-rata terdapat satu peristiwa setiap 10
menit (lihat Contoh \ref{exam 2.2.7.5}). Anda menjumpai sistem ini pada waktu
100 dan menunggu hingga peristiwa berikutnya. Buatlah dugaan tentang rata-rata
lama Anda harus menunggu. Tulislah program untuk memeriksa apakah dugaan Anda
benar.

\i\label{exer 2.2.8.6} Seperti dalam Latihan \ref{exer 2.2.8.5}, anggaplah
kita mempunyai suatu barisan peristiwa, tetapi sekarang anggaplah bahwa waktu
$X$ antara dua peristiwa berdistribusi seragam antara 5 dan 15. Seperti
sebelumnya, Anda menjumpai sistem ini pada waktu 100 dan menunggu hingga
peristiwa berikutnya. Buatlah dugaan tentang rata-rata lama Anda harus
menunggu. Tulislah program untuk memeriksa apakah dugaan Anda benar.

\i\label{exer 2.2.11} Untuk contoh seperti dalam Latihan
\ref{exer 2.2.8.5} dan \ref{exer 2.2.8.6}, mungkin tampak bahwa setidaknya Anda
seharusnya tidak perlu menunggu rata-rata \emx {lebih} dari 10 menit jika
rata-rata waktu antara dua peristiwa adalah 10 menit. Sayangnya, bahkan hal ini
tidak benar. Untuk melihat alasannya, tinjaulah asumsi berikut tentang waktu
antara dua peristiwa. Anggaplah waktu antara dua peristiwa adalah 3 menit
dengan peluang .9 dan 73 menit dengan peluang .1. Tunjukkan melalui simulasi
bahwa rata-rata waktu antara dua peristiwa adalah 10 menit, tetapi jika Anda
menjumpai sistem ini pada waktu 100, rata-rata waktu tunggu Anda lebih dari 10
menit.

\i\label{exer 2.2.13} Ambillah sebuah tongkat
\index{tongkat sepanjang satu satuan} sepanjang satu satuan dan patahkan menjadi
tiga bagian dengan memilih titik patahnya secara acak. (Titik-titik patah
dianggap dipilih secara bersamaan.) Berapakah peluang bahwa ketiga bagian dapat
digunakan untuk membentuk segitiga? \emx {Petunjuk}: jumlah panjang sebarang dua
bagian harus melebihi panjang bagian ketiga, sehingga setiap bagian harus
memiliki panjang $< 1/2$. Sekarang gunakan Latihan~\ref{exer 2.2.8}(g).

\i\label{exer 2.2.14} Ambillah sebuah tongkat sepanjang satu satuan
\index{tongkat sepanjang satu satuan} dan patahkan menjadi dua bagian dengan
memilih titik patahnya secara acak. Kemudian patahkan bagian yang lebih panjang
pada suatu titik acak. Berapakah peluang bahwa ketiga bagian dapat digunakan
untuk membentuk segitiga?

\i\label{ex:cr} Pilihlah secara saling bebas dua bilangan $B$ dan $C$
\emx {secara acak} dari interval $[-1,1]$ dengan distribusi seragam, lalu
tinjaulah persamaan kuadrat\index{persamaan kuadrat, akar}
$$
x^2 + Bx + C = 0\ .
$$
Carilah peluang bahwa akar-akar persamaan ini
\begin{enumerate}
\item keduanya real.

\item keduanya positif.
\end{enumerate}

\emx {Petunjuk}: (a) mensyaratkan $0 \leq B^2 - 4C$,
(b) mensyaratkan $0 \leq B^2 - 4C$, $B \leq 0$, $0 \leq C$.

\i\label{ex:cs} Di Tunbridge World's Fair, suatu permainan lempar koin
berlangsung sebagai berikut. Koin seperempat dolar dilempar ke papan catur.
Pengelola menyimpan semua koin tersebut, tetapi untuk setiap koin yang jatuh
sepenuhnya di dalam satu petak papan catur, pengelola membayar satu dolar.
Anggaplah sisi setiap petak dua kali diameter koin seperempat dolar dan hasilnya
dinyatakan oleh koordinat yang dipilih \emx {secara acak.} Apakah permainan ini
adil?

\i\label{exer 2.2.16} Tiga titik dipilih \emx {secara acak} pada lingkaran
dengan \emx {keliling satu satuan.} Berapakah peluang bahwa segitiga
\index{segitiga!lancip} yang titik-titik sudutnya adalah ketiga titik tersebut
mempunyai tiga sudut lancip? \emx {Petunjuk}: salah satu sudut tumpul jika dan
hanya jika ketiga titik terletak pada setengah lingkaran yang sama. Pandanglah
keliling sebagai interval $[0,1]$. Tempatkan satu titik di 0 dan titik-titik
lainnya di $B$ dan $C$.

\i\label{exer 2.2.17} Tulislah program untuk memilih bilangan acak $X$ dalam
interval $[2,10]$ sebanyak 1000 kali dan mencatat proporsi hasil yang memenuhi
$X > 5$, proporsi yang memenuhi $5 < X < 7$, dan proporsi yang memenuhi
$x^2 - 12x + 35 > 0$. Bagaimana hasil ini dibandingkan dengan Latihan
\ref{exer 2.2.1}?

\i\label{exer 2.2.18} Tulislah program untuk memilih suatu titik $(X,Y)$
\emx {secara acak} dalam persegi dengan panjang sisi 20 inci, lakukan ini
10{,}000 kali, dan catat proporsi hasil yang terletak dalam jarak 19 inci dari
pusat; di antara hasil tersebut, proporsi yang terletak antara 8 dan 10 inci
dari pusat; dan, di antara hasil tersebut, proporsi yang terletak dalam kuadran
pertama persegi. Bagaimana hasil ini dibandingkan dengan hasil
Latihan~\ref{exer 2.2.4}?

\i\label{exer 2.2.19} Tulislah program untuk menyimulasikan masalah yang
diuraikan dalam Latihan~\ref{exer 2.2.7} (lihat
Latihan~\ref{exer 2.2.17}). Bagaimana hasil simulasi dibandingkan dengan hasil
Latihan~\ref{exer 2.2.7}?

\i\label{exer 2.2.20} Tulislah program untuk menyimulasikan masalah yang
diuraikan dalam Latihan \ref{exer
2.2.13}.

\i\label{exer 2.2.21} Tulislah program untuk menyimulasikan masalah yang
diuraikan dalam Latihan \ref{exer
2.2.16}.

\i\label{exer 2.2.22} Tulislah program untuk menjalankan percobaan berikut.
Sebuah koin dilempar 100 kali dan jumlah kepala yang muncul dicatat. Percobaan
ini kemudian diulang 1000 kali. Buatlah program Anda menggambar diagram batang
untuk proporsi dari 1000 percobaan ketika jumlah kepala adalah $n$, untuk setiap
$n$ dalam interval $[35,65]$. Apakah diagram batang tersebut tampak dapat
disesuaikan dengan kurva normal?

\item\label{exer 2.2.23} Tulislah program yang memilih bilangan acak antara 0
dan 1 lalu menghitung negatif dari logaritmanya. Ulangi proses ini berkali-kali
dan gambarlah diagram batang yang menunjukkan berapa kali hasilnya jatuh dalam
setiap interval sepanjang 0.1 dalam $[0,10]$. Pada diagram batang ini,
gambarlah grafik fungsi kepadatan $f(x) = e^{-x}$. Seberapa baik fungsi
kepadatan ini sesuai dengan diagram Anda?

\end{LJSItem}
%\end{document}

